\documentclass[10pt,a4paper]{amsart}
\usepackage[margin=19mm]{geometry}
\usepackage{amsmath,amssymb,mathtools}
\usepackage[scr=rsfs,cal=euler]{mathalfa}
\usepackage{xfrac}
\usepackage[unicode,hidelinks,bookmarks=false]{hyperref}
\newcommand{\ie}{i.e.\ }
\newcommand{\cf}{cf.\ }
\newcommand{\resp}{respectively\ }
\newcommand{\Subsection}{\subsection}
\providecommand{\nobreakdash}{\nobreak\hbox{-}}
\setlength{\emergencystretch}{2em}
\multlinegap0pt
\usepackage{amssymb}
\newtheorem{theorem}{Theorem}
\title{Cyclic $p$-$\varphi$-contraction mappings}
\author{Seyyed Mohammad Sadegh Nabavi Sales}
\date{Source: arXiv:2602.16226v1 (CC BY 4.0)}
\begin{document}
\maketitle

\noindent\emph{Source and licence.} Cyclic $p$-$\varphi$-contraction mappings. Version arXiv:2602.16226v1 (CC BY 4.0), distributed by arXiv under the Creative Commons Attribution 4.0 International licence, http://creativecommons.org/licenses/by/4.0/, as recorded in the arXiv licence metadata for this exact version and shown on the abstract page https://arxiv.org/abs/2602.16226v1. Full licence notice: \texttt{licenses/arxiv-2602.16226-CC-BY-4.0.txt}.\par\smallskip

\noindent\textit{Editorial scope.} Counted unit: the article's theorem SS1 (source0.tex lines 109--111). The article's complete original proof of it is delivered with it (source0.tex lines 112--154, one \texttt{proof} environment of 43 lines). No mathematics is written, repaired or re-derived: the statement and the proof are the article's own lines, in the article's own notation, and no retained line is substituted or re-flowed.  No further source-local context is needed: every cross-reference of every retained line resolves inside the two delivered spans.  Nothing was omitted from any retained span, so the package carries no omission record.  No retained line contains a citation, so the wrapper carries no bibliography and no bibliography entry.  No retained line contains a figure environment, an image-inclusion command or drawing code, so no image is delivered and the package declares no asset dependency.  Editorial formatting only, in addition: an AMS \texttt{amsart} preamble that restates the article's own declarations and macros used by the retained lines, namely the environments \texttt{theorem} on separate counters, together with the standard packages those lines need.  The retained environments print in this document as Theorem 1 in order of appearance, whereas the article prints the same items with the numbering of its own full document.  The complete native source of the article as distributed in the arXiv e-print for this exact version is bundled unchanged as \texttt{sources/194860/source0.tex}.
\begin{theorem}\label{SS1}
 If $\{x_{mn}\}_{n\geq0}$ has a convergent subsequence in $A_1$, then there exists $x\in A_1$ so that $$d_p(x,Tx,\cdots,T^{m-1}x)=d_p(A_1,A_2,\cdots,A_m).$$
\end{theorem}

\begin{proof}
We may assume that  $\{x_{mn}\}_{n\geq0}$ is convergent and $x_{mn}\to x$ as $n\to\infty$. First assume that $p=\infty$. Therefore,
\begin{eqnarray*}
d_\infty(A_1,\cdots,A_m)&\leq& d_\infty(x,x_{mn+1},x_{mn+2},\cdots,x_{mn+m-1})\\
&=&\max\{d(x,x_{mn+1}),d(x_{mn+1},x_{mn+2}),\cdots,d(x_{mn+m-1},x)\}\\
&\leq&\max\{d(x,x_{mn})+d(x_{mn},x_{mn+1}),d(x,x_{mn})+d(x_{mn+1},x_{mn+2})\\&&,\cdots,d(x,x_{mn})+d(x_{mn+m-1},x_{mn})\}\\&=&d(x,x_{mn})+d_\infty(x_{mn},x_{mn+1},\cdots,x_{mn+m-1})\to d_\infty(A_1,\cdots,A_m).
\end{eqnarray*}
It follows that
\begin{eqnarray*}
d_\infty(A_1,\cdots,A_m)&\leq& d_\infty(x_{mn},Tx,x_{mn+2},\cdots,x_{mn+m-1})\\
&\leq&d_\infty(x_{mn-1},x,x_{mn+1}\cdots,x_{mn+m-2})\to d_\infty(A_1,\cdots,A_m).
\end{eqnarray*}
Hence $d_\infty(x_{mn},Tx,x_{mn+2},\cdots,x_{mn+m-1})\to d_\infty(A_1,\cdots,A_m)$ which implies that
\begin{eqnarray*}
d_\infty(A_1,\cdots,A_m)&\leq& d_\infty(x,Tx,x_{mn+2},\cdots,x_{mn+m-1})\\
&\leq&d(x,x_{mn})+d_\infty(x_{mn},Tx,x_{mn+1}\cdots,x_{mn+m-1})\\&&\to d_\infty(A_1,\cdots,A_m)
\end{eqnarray*}
that means $d_\infty(x,Tx,x_{mn+2},\cdots,x_{mn+m-1})\to d_\infty(A_1,\cdots,A_m)$. So, by an easy induction we can show that
$$d_\infty(x_{mn},Tx,T^2x,\cdots,T^{m-2}x,T^{m-1}x)\to d_\infty(A_1,\cdots,A_m).$$
This yeilds that
\begin{eqnarray*}
d_\infty(A_1,\cdots,A_m)&\leq&d_\infty(x,Tx,T^2x,\cdots,T^{m-2}x,T^{m-1}x)\\
&\leq&d(x,x_{mn})+d_\infty(x_{mn},Tx,T^2x,\cdots,T^{m-2}x,T^{m-1}x)\\&&\to d_\infty(A_1,\cdots,A_m).
\end{eqnarray*}
Now, assume that $1\leq p<\infty$. We have that
\begin{eqnarray*}
d_p(A_1,\cdots,A_m)&\leq& d_p(x,x_{mn+1},x_{mn+2},\cdots,x_{mn+m-1})\\
&=&\{d(x,x_{mn+1})^p+d(x_{mn+1},x_{mn+2})^p+\cdots+d(x_{mn+m-1},x)^p\}^{1\over p}\\
&\leq&\{(d(x,x_{mn})+d(x_{mn},x_{mn+1}))^p+d(x_{mn+1},x_{mn+2})^p\\&&+\cdots+(d(x,x_{mn})+d(x_{mn+m-1},x_{mn}))^p\}^{1\over p}\\&\leq&2^{1\over p}d(x,x_{mn})+d_p(x_{mn},x_{mn+1},\cdots,x_{mn+m-1})\to d_p(A_1,\cdots,A_m),
\end{eqnarray*}
which implies that $d_p(x,x_{mn+1},x_{mn+2},\cdots,x_{mn+m-1})\to d_p(A_1,\cdots,A_m)$.
Thus
\begin{eqnarray*}
d_p(A_1,\cdots,A_m)&\leq& d_p(x,Tx,x_{mn+2},\cdots,x_{mn+m-1})\\
&\leq&2^{1\over p}d(x,x_{mn})+d_p(x_{mn},Tx,x_{mn+1}\cdots,x_{mn+m-2})\\&&\to d_p(A_1,\cdots,A_m),
\end{eqnarray*}
that means $ d_p(x,Tx,x_{mn+2},\cdots,x_{mn+m-1})\to d_p(A_1,\cdots,A_m).$ An easy induction shows that $$d_p(x_{mn},Tx,T^2x,\cdots,T^{m-2}x,T^{m-1}x)\to d_p(A_1,\cdots,A_m)$$
and this yeilds that
\begin{eqnarray*}
d_p(A_1,\cdots,A_m)&\leq&d_p(x,Tx,T^2x,\cdots,T^{m-2}x,T^{m-1}x)\\
&\leq&2^{1\over p}d(x,x_{mn})+d_p(x_{mn},Tx,T^2x,\cdots,T^{m-2}x,T^{m-1}x)\\&&\to d_p(A_1,\cdots,A_m).
\end{eqnarray*}
\end{proof}

\end{document}
